\documentclass[10pt,a4paper]{amsart}
\usepackage[margin=19mm]{geometry}
\usepackage{amsmath,amssymb,mathtools}
\usepackage[scr=rsfs,cal=euler]{mathalfa}
\usepackage{xfrac}
\usepackage[unicode,hidelinks,bookmarks=false]{hyperref}
\newcommand{\ie}{i.e.\ }
\newcommand{\cf}{cf.\ }
\newcommand{\resp}{respectively\ }
\newcommand{\Subsection}{\subsection}
\providecommand{\nobreakdash}{\nobreak\hbox{-}}
\setlength{\emergencystretch}{2em}
\multlinegap0pt
\usepackage{amssymb}
\usepackage[shortlabels]{enumitem}
\usepackage{xcolor}
\usepackage{tikz}
\newtheorem{lemma}{Lemma}
\renewcommand{\phi}{\varphi}
\title{Spherical subgroups in reductive algebraic groups}
\author{Friedrich Knop, Gerhard Röhrle}
\date{Source: arXiv:2607.03864v2 (CC BY 4.0)}
\begin{document}
\maketitle

\noindent\emph{Source and licence.} Spherical subgroups in reductive algebraic groups. Version arXiv:2607.03864v2 (CC BY 4.0), distributed by arXiv under the Creative Commons Attribution 4.0 International licence, http://creativecommons.org/licenses/by/4.0/, as recorded in the arXiv licence metadata for this exact version and shown on the abstract page https://arxiv.org/abs/2607.03864v2. Full licence notice: \texttt{licenses/arxiv-2607.03864-CC-BY-4.0.txt}.\par\smallskip

\noindent\textit{Editorial scope.} Counted unit: the article's lemma lemma:hard (source0.tex lines 609--614). The article's complete original proof of it is delivered with it (source0.tex lines 616--626, one \texttt{proof} environment of 11 lines). No mathematics is written, repaired or re-derived: the statement and the proof are the article's own lines, in the article's own notation, and no retained line is substituted or re-flowed.  No further source-local context is needed: every cross-reference of every retained line resolves inside the two delivered spans.  Nothing was omitted from any retained span, so the package carries no omission record.  No retained line contains a citation, so the wrapper carries no bibliography and no bibliography entry.  No retained line contains a figure environment, an image-inclusion command or drawing code, so no image is delivered and the package declares no asset dependency.  Editorial formatting only, in addition: an AMS \texttt{amsart} preamble that restates the article's own declarations and macros used by the retained lines, namely the environments \texttt{lemma} on separate counters, together with the standard packages those lines need.  The retained environments print in this document as Lemma 1 in order of appearance, whereas the article prints the same items with the numbering of its own full document.  The complete native source of the article as distributed in the arXiv e-print for this exact version is bundled unchanged as \texttt{sources/204429/source0.tex}.
\begin{lemma}[Strong Transitivity]\label{lemma:hard}
  Let $G$ be a connected, reductive group and let
  $H_1\subseteq H_2\subseteq G$ be connected, reductive subgroups. Put
  $L:=L(G/H_2)$. Then $G/H_1$ is spherical if and only if $G/H_2$ is
  spherical and $H_2/H_1$ is spherical as an $L$-variety.
\end{lemma}

\begin{proof}
  Consider the canonical morphism $\phi:X_1:=G/H_1\to X_2:=G/H_2$. Let $x$ be a
  point in the open $B$-orbit $X_2^0$ of $X_2$. Since $X_2^0$ is also
  $P$-invariant, we get a bijective morphism $P/C_P(x)\to X_2^0$, hence a
  bijective morphism
  $P\times^{C_P(x)}F\to\phi^{-1}(X_2^0)\subseteq X_1$ with
  $F :=\phi^{-1}(x)\cong H_2/H_1$. Thus, $B$ has an open orbit in $X_1$
  if and only if $C_P(x)=C_B(x)$ has an open orbit in $F$. Now
  $C_B(x)$ is a Borel subgroup of $C_P(x)$. Thus $X_1$ is spherical for $G$ if
  and only if $F$ is spherical as an $L (=C_P(x)^0)$-variety.
\end{proof}

\end{document}
